\uuid{5bQE}
\titre{Division euclidienne}
\auteur{bodin ; hanani ; mzari}
\organisation{exo7}
\licence{CC BY-NC-SA 4.0}
\datecreate{2020-01-20}
\niveau{L1}
\module{Algèbre}
\chapitre{Polynôme, fraction rationnelle}
\sousChapitre{Division euclidienne}
\difficulte{3}
\theme{Arithmétique des polynômes, division euclidienne}
\exo7qcmid{lille-1.polynomes.015}

\contenu{
\begin{question}
\qclassification{105.01,105.02}
\qtags{motcle=division euclidienne}
Soit $A$ un polynôme de degré $n\ge1$. Soit $B$ un polynôme de degré $m\ge1$, avec $m \le n$.
Soit $A = B \times Q + R$ la division euclidienne de $A$ par $B$. On note
$q = \deg Q$ et $r = \deg R$ (avec $r=-\infty$ si $R=0$).
Quelles sont les assertions vraies (quelque soient $A$ et $B$) ?
\begin{answers}
    \good{$q = n-m$}

    \good{$r < m$}

    \bad{$r=0 \implies A$ divise $B$.}

    \bad{$n = mq + r$}
\end{answers}
\begin{explanations}
On a $\deg R < \deg B$. Il ne faut pas confondre $R=0$ et $r=0$.
En plus $\deg(A) = \deg(B\times Q) = \deg(A) + \deg(Q)$.
\end{explanations}
\end{question}
}
